% =============================================================================
% Chapitre : Méthodes à noyaux et SVM
% =============================================================================

\minitoc

L'objectif des méthodes à noyaux est de transformer une mesure heuristique de similarité entre données en un algorithme d'apprentissage non linéaire. En s'appuyant implicitement sur un espace de Hilbert de dimension arbitraire (potentiellement infinie), on conserve toute la simplicité de l'optimisation linéaire tout en évitant d'avoir à calculer explicitement les plongements.

% ======================================================================
% Motivation et Théorème du Représentant

\section{De la similarité au Théorème du Représentant}

% --------------------------------------------
\subsection{Intuition géométrique et plongement}

La plupart des jeux de données ne sont pas linéairement séparables dans leur espace d'entrée d'origine $\mathcal{X}$. L'idée fondamentale consiste à envoyer les données dans un espace de caractéristiques de plus grande dimension (un espace de Hilbert $\mathcal{H}$) via une \emph{feature map} :
\begin{equation}
    \varphi : \mathcal{X} \longrightarrow \mathcal{H}
\end{equation}
Dans $\mathcal{H}$, un hyperplan linéaire peut alors séparer des classes qui formaient des régions géométriques complexes ou intriquées dans $\mathcal{X}$. 
On formule alors le problème sous la forme d'une minimisation de risque empirique régularisé sur la totalité de l'espace hilbertien $\mathcal{H}$ :
\begin{equation}
    \hat{\theta}_S \in \arg\min_{\theta \in \mathcal{H}} \; \frac{1}{n} \sum_{i=1}^n l\left(y_i, \langle \varphi(x_i), \theta \rangle_{\mathcal{H}}\right) + \frac{\lambda}{2} \|\theta\|_{\mathcal{H}}^2 \label{eq:initial-problem}
\end{equation}
où $\lambda > 0$ contrôle la régularisation. 

\begin{remark}[Obstruction numérique]
    Si $\mathcal{H}$ est de dimension infinie, on ne peut ni représenter le vecteur de paramètres $\theta \in \mathcal{H}$, ni évaluer explicitement $\varphi(x)$ par descente de gradient standard. Le théorème suivant résout ce paradoxe.
\end{remark}

% --------------------------------------------
\subsection{Le Théorème du Représentant}

\begin{theorem}[Théorème du Représentant]
    Soit $\mathcal{X}$ un ensemble, $\mathcal{H}$ un espace de Hilbert réel muni du produit scalaire $\langle \cdot, \cdot \rangle_{\mathcal{H}}$ et $\varphi : \mathcal{X} \to \mathcal{H}$. Soit un échantillon $S = (x_1, \dots, x_n)$ et le sous-espace engendré par les observations :
    \begin{equation}
        \mathcal{H}_D := \operatorname{vect}\{\varphi(x_1), \dots, \varphi(x_n)\} \subset \mathcal{H}
    \end{equation}
    Soit $\Psi : \R^n \times \R_+ \to \R$ une fonction croissante par rapport à son second argument. En notant $u(\theta) = (\langle \varphi(x_1), \theta \rangle_{\mathcal{H}}, \dots, \langle \varphi(x_n), \theta \rangle_{\mathcal{H}})^\top \in \R^n$, on a :
    \begin{equation}
        \inf_{\theta \in \mathcal{H}} \Psi\left(u(\theta), \|\theta\|_{\mathcal{H}}^2\right) = \inf_{\theta \in \mathcal{H}_D} \Psi\left(u(\theta), \|\theta\|_{\mathcal{H}}^2\right)
    \end{equation}
    Si l'infimum est atteint, il l'est en un point $\hat{\theta}_S \in \mathcal{H}_D$. De plus, si $\Psi$ est \emph{strictement croissante} par rapport à son second argument, tout minimiseur appartient nécessairement à $\mathcal{H}_D$.
\end{theorem}

% TODO : Rajouter preuve du théorème par Pythagore 

Grâce à ce théorème, on cherche donc $ \hat{ \theta}_S \in \mathcal{H}_D$ de dimension finie. Puisque l'on en connaît une famille génératrice, on cherche donc : 
$$ \alpha = ( \alpha_1, \dots, \alpha_n )^\top \in \R^n, \quad \text{tels que} \quad \hat{ \theta}_S = \sum_{i=1}^{n} \alpha_i \varphi(x_i). $$
Ainsi, la dimension du problème d'optimisation ne dépend plus de celle de $ \mathcal{H}$ mais seulement du nombre de d'échantillons $n$. Déterminons maintenant comment trouver ce $ \hat{ \alpha}_S$. 

\begin{proposition}[Kernel Trick]
    Remplaçons $ \theta$ par sa nouvelle expression dans le problème de départ \ref{eq:initial-problem} : 
    \begin{itemize}
        \item La norme devient : 
        \begin{align*}
            \| \theta \|^2_{ \mathcal{H}} &= \Bigl\langle \sum_{i=1}^{n} \alpha_i \varphi(x_i), \sum_{j=1}^{n} \alpha_j \varphi(x_j) \Bigr\rangle_{ \mathcal{H}}
            = \sum_{i,j = 1}^{n} \alpha_i \alpha_j \bigl\langle \varphi(x_i), \varphi(x_j) \bigr\rangle_{ \mathcal{H}} \\ 
            &=  \sum_{i,j = 1}^{n} \alpha_i \alpha_j K_{i,j} \quad \text{avec} \quad K = \left( \bigl\langle \varphi(x_i), \varphi(x_j) \bigr\rangle_{ \mathcal{H}}\right)_{i,j}
        \end{align*}
        \item Les prédictions deviennent : 
        \begin{align*}
            \langle \varphi(x_i), \theta \rangle_{ \mathcal{H}} = \sum_{j=1}^{n} \alpha_j \langle \varphi(x_i), \varphi(x_j) \rangle_{ \mathcal{H}} = (K \alpha)_i
        \end{align*}
    \end{itemize}
    On a donc le problème réécrit suivant : 
    \begin{equation}
        \hat{ \alpha}_S \in \underset{ \alpha \in \R^n}{\arg \min} \frac{1}{n} \sum_{i=1}^{n} l \left(y_i , (K \alpha)_i\right) + \frac{ \lambda}{2} \alpha^\top K \alpha. 
    \end{equation}
    On remarque que l'on a plus besoin de calculer explicitement $ \varphi$ ni $f$ : tout dépend de ce que l'on appelera le \emph{noyau} $K$. 
    Notons ici $K = (k(x_i, x_j))_{i,j}$ tel que $k(x, x_j) = \langle \varphi(x), \varphi(x_j) \rangle_{ \mathcal{H}}$. Ainsi, pour l'entraînement nous n'avons besoin que de la matrice de Gram $K$ de taille $ n \times n$. Pour l'inférence sur un $x$ donné, nous notons $ k_x = (k(x, x_1), \dots, k(x, x_n)) \in \R^n$ tel que $ f_{ \hat{ \theta}_S} = \hat{ \alpha}_S^\top k_x$. 
\end{proposition}

On peut maintenant remarquer que la \emph{feature map} $ \varphi$ n'est jamais calculée et que les prédictions en $x$ utilisent seulement les $n$ valeurs $ \langle \varphi(x), \varphi(x_i) \rangle_{ \mathcal{H}}$. Ainsi, au lieu de définir un espace $ \mathcal{H}$ dans lequel nous ne pouvons faire aucun calcul, on pourrait directement choisir $k$ telle que $ k : \mathcal{X} \times \mathcal{X} \longrightarrow \R$ de la forme $ k(x, x') = \langle \varphi(x), \varphi(x') \rangle$. C'est le sujet de la section suivante qui définit les \emph{noyaux} et introduit la théorie nécessaire dessus. 

% ======================================================================
% Noyaux Définis Positifs et RKHS

\section{Noyaux définis positifs, Théorème de Bochner, et Théorème d'Aronszajn}

Dans la partie précédente, nous avons reformulé le problème de minimisation du risque empirique à l'aide d'un plongement dans un espace de Hilbert $ \mathcal{H}$. Nous avons ensuite remarqué que tous les calculs dans cet espace dépendent en réalité d'une \emph{fonction de similitude} encodée dans un \emph{noyau} $K$. 

Or, en pratique, on raisonne dans le sens inverse : à partir du jeu de données, on définit un noyau $K$ et on souhaite savoir si l'on peut effectivement utiliser ce noyau pour notre minimisation. C'est le théorème d'Aronszajn qui nous l'assure. Mais avant de l'aborder, définissons formellement la notion de noyau ainsi que ses propriétés. 

% --------------------------------------------
\subsection{Noyau SDP : Définition et propriétés}

Au lieu de définir $\mathcal{H}$ puis $\varphi$, on choisit directement une fonction $k : \mathcal{X} \times \mathcal{X} \to \R$ et on cherche à quelles conditions elle correspond à un produit scalaire hilbertien (i.e un noyau semi-défini positif). 

\begin{definition}[Noyau semi-défini positif]
    Une application $k : \mathcal{X} \times \mathcal{X} \to \R$ est appelée un \emph{noyau semi-défini positif} (noyau SDP ou PDSK) si elle vérifie :
    \begin{enumerate}
        \item \textbf{Symétrie} : $\forall x, x' \in \mathcal{X}, \quad k(x, x') = k(x', x)$.
        \item \textbf{Positivité} : $\forall n \geqslant 1, \; \forall (x_1, \dots, x_n) \in \mathcal{X}^n, \; \forall (\alpha_1, \dots, \alpha_n) \in \R^n$ :
        \begin{equation}
            \sum_{i,j=1}^n \alpha_i \alpha_j k(x_i, x_j) \geqslant 0 \iff K \in \mathcal{S}_n^+(\R)
        \end{equation}
    \end{enumerate}
\end{definition}

\begin{prop}[Caractérisation par matrice de Gram]
    La matrice de Gram $ G = \left(\bigl\langle u_i, u_j\bigr\rangle\right)_{i,j}$ associée à toute famille finie $ (u_1, \dots, u_n)$ est symétrique semi-définie positive.
\end{prop}

\begin{proposition}[Tout produit scalaire définit un noyau SDP]
    Si $\mathcal{H}$ est un espace de Hilbert et $\varphi : \mathcal{X} \to \mathcal{H}$, alors $k(x, x') := \langle \varphi(x), \varphi(x') \rangle_{\mathcal{H}}$ est un noyau SDP.
\end{proposition}

\begin{example}
    Voici quelques exemples de noyaux SDP : 
    \begin{itemize}
        \item Noyau \emph{linéaire} : $ k(x, x') = x^\top x'$ sur $ \mathcal{X} = \R^d$. 
        \item Noyau \emph{polynomial} : $ k(x, x') = (x^\top x')^q$ pour $ q \in \N^*$. 
        \item Noyau \emph{exponentiel} : $ k(x, x') = \exp \left(- a \| x - x'\|_2^2 \right)$ pour $ a > 0$. 
    \end{itemize}
\end{example}

Les noyaux SDP peuvent être construit par opérations. Le théorème suivant détaille les opérations possibles entre noyaux SDP. 

\begin{theorem}[Opérations]
    L'ensemble des noyaux semi-définis positifs sur un espace $ \mathcal{X}$ est stable par : 
    \begin{itemize}
        \item multiplication scalaire par un coefficient positif ou nul, 
        \item somme finie, 
        \item produits, 
        \item limite ponctuelle quand elle existe, 
        \item conjugaison par une fonction. 
    \end{itemize}
\end{theorem}

Grâce à ce théorème, nous pouvons montrer que beaucoup d'application $k$ sont des noyaux. Voyons-en quelques unes. 

\begin{proposition}
    Soit $ q \geqslant 1$ un entier, alors $(x, x') \longmapsto (x^\top x')^q$ et $(x, x') \longmapsto (1 + x^\top x')^q$ sont des noyaux SDP. 
\end{proposition}

\begin{proof}
    Pour le premier, $(x, x') \mapsto x^\top x'$ est un produit scalaire donc c'est un noyau. À la puissance $q$, c'est une multiplication finie. C'est donc un noyau par induction sur $q$. 

    Pour le second, on identifie $ 1 + x^\top x' = k(x, x') + x^\top x'$ avec $ k(x, x') = 1$. Puisque $ \sum_{i,j} \alpha_i \alpha_j = \left( \sum_{i} \alpha_i \right)^2 \geqslant 0$, alors c'est un noyau. Par somme et d'après le résultat précédent, $(x, x') \longmapsto (1 + x^\top x')^q$ est aussi un noyau SDP. 
\end{proof}

\begin{proposition}
    Soit $k$ un noyau SDP, alors $ \exp(k) : (x, x') \longmapsto \exp(k(x, x'))$ est un noyau SDP. 
\end{proposition}

\begin{proof}
    Si $k$ est un noyau alors, $ \forall x, x' \in \mathcal{X}$, on a : 
    $$ \exp(k(x, x')) = \sum_{n=0}^{\infty} \frac{k(x, x')^n}{n !}. $$
    Or, pour tout $ n \geqslant 0$, $\frac{k(x, x')^n}{n !}$ est un noyau par produit fini et multiplication par un scalaire positif. Par somme, $\sum_{i=1}^{n} \frac{k(x, x')^i}{i !}$ en est un aussi. Enfin, par définition de l'exponentielle, $ \exp a $ converge pour tout $a \in \R$, donc par limite de noyaux SDP, $ \exp (k)$ est aussi un noyau SDP. 
\end{proof}

% --------------------------------------------
\subsection{Construire des noyaux SDP : Théorème de Bochner}

Montrer qu'un noyau est SDP est simple lorsqu'il dérive de produits scalaires combinés algébriquement. En revanche, si l'on cherche à mesurer la proximité géométrique locale entre deux observations, le calcul direct devient inextricable. Par exemple, pour le noyau gaussien $k(x, x') = \exp(-a \|x - x'\|_2^2)$, vérifier point par point que :
\begin{equation*}
    \forall n \geqslant 1, \quad \forall (x_1, \dots, x_n) \in (\R^d)^n, \quad \forall (\alpha_1, \dots, \alpha_n) \in \R^n, \quad \sum_{i=1}^n \sum_{j=1}^n \alpha_i \alpha_j \exp(-a \|x_i - x_j\|_2^2) \geqslant 0
\end{equation*}
est difficile sans outil analytique. Le théorème de Bochner permet de lever ce verrou en transposant la question de l'espace spatial à l'espace fréquentiel via la transformée de Fourier.

\begin{theorem}[Théorème de Bochner -- Invariance par translation]
    Soit $q : \R^d \to \R$ continue. Le noyau stationnaire $k(x, x') := q(x - x')$ est semi-défini positif si et seulement si $q$ est la transformée de Fourier inverse d'une mesure borélienne positive finie $\rho$ sur $\R^d$ :
    \begin{equation}
        q(z) = \frac{1}{(2\pi)^d} \int_{\R^d} e^{i z^\top \omega} \, \mathrm{d}\rho(\omega).
    \end{equation}
\end{theorem}

\begin{proof}[Preuve de la condition suffisante]
    Supposons que $q$ s'écrive sous forme d'intégrale de Fourier par rapport à une mesure positive finie $\rho$. Pour toute famille $(x_1, \dots, x_n)$ et tout vecteur $\alpha \in \R^n$ :
    \begin{align*}
        \sum_{i=1}^n \sum_{j=1}^n \alpha_i \alpha_j q(x_i - x_j) 
        &= \sum_{i,j=1}^n \alpha_i \alpha_j \frac{1}{(2 \pi)^d} \int_{\R^d} e^{i (x_i - x_j)^\top \omega} \, \mathrm{d}\rho(\omega) \\ 
        &= \frac{1}{(2 \pi)^d} \int_{\R^d} \sum_{i=1}^n \alpha_i e^{i x_i^\top \omega} \sum_{j=1}^n \alpha_j e^{-i x_j^\top \omega} \, \mathrm{d}\rho(\omega) \\ 
        &= \frac{1}{(2 \pi)^d} \int_{\R^d} \left( \sum_{i=1}^n \alpha_i e^{i x_i^\top \omega} \right) \overline{\left( \sum_{j=1}^n \alpha_j e^{i x_j^\top \omega} \right)} \, \mathrm{d}\rho(\omega) \\
        &= \frac{1}{(2 \pi)^d} \int_{\R^d} \underbrace{\left| \sum_{i=1}^n \alpha_i e^{i x_i^\top \omega} \right|^2}_{\geqslant 0} \, \underbrace{\mathrm{d}\rho(\omega)}_{\geqslant 0} \geqslant 0.
    \end{align*}
    La fonction $q$ étant paire dès lors que la mesure est symétrique ($\rho(-\omega) = \rho(\omega)$), le noyau est bien symétrique et semi-défini positif.
\end{proof}

\begin{example}[Cas du noyau Gaussien / RBF]
    Pour $a > 0$ et $z \in \R^d$, la fonction gaussienne $q(z) = \exp(-a \|z\|_2^2)$ admet pour transformée de Fourier inverse :
    \begin{equation*}
        e^{-a \|z\|_2^2} = \frac{1}{(2\pi)^d} \int_{\R^d} e^{i z^\top \omega} \left( \frac{\pi}{a} \right)^{d/2} e^{-\frac{\|\omega\|_2^2}{4a}} \, \mathrm{d}\omega.
    \end{equation*}
    La mesure $\mathrm{d}\rho(\omega) = (\pi/a)^{d/2} e^{-\|\omega\|_2^2/(4a)} \, \mathrm{d}\omega$ est une mesure gaussienne positive, finie et invariante par $\omega \mapsto -\omega$. Le théorème de Bochner garantit immédiatement que le noyau gaussien est semi-défini positif.
\end{example}

% --------------------------------------------
\subsection{Théorème d'Aronszajn}

Reprenons le raisonnement de ce chapitre. Le théorème du représentant nous a montré que pour notre problème d'optimisation, si on a un espace de Hilbert $ \mathcal{H}$ et un application $ \varphi$, il nous suffit seulement de calculer des produits scalaires $ k(x, x') = \langle \varphi(x), \varphi(x') \rangle_{ \mathcal{H}}$. 

Ensuite, nous avons développé une théorie sur les noyaux SDP pour les construire facilement. Le théorème de Bochner nous permet d'en construire autant que l'on souhaite à partir de mesures positives finies. 

Or, nous ne sommes pas assurés que pour un tel noyau construit spécifiquement pour nos données, il existe un espace $ \mathcal{H}$ et un plongement $ \varphi$ qui convienne. C'est ici que le théorème d'Aronszajn intervient : il nous assure que tout noyau SDP est le produit scalaire d'un espace unique de fonctions appelé RKHS. 

\begin{theorem}[Théorème d'Aronszajn]
    Soit $k : \mathcal{X} \times \mathcal{X} \to \R$ un noyau symétrique semi-défini positif. Alors, il existe un unique espace de Hilbert $\mathcal{H}_k$ constitué de fonctions de $\mathcal{X}$ dans $\R$, et une application canonique $\varphi_k : \mathcal{X} \to \mathcal{H}_k$, tels que :
    \begin{enumerate}
        \item Pour tout $x \in \mathcal{X}$, la fonction $k_x := k(x, \cdot)$ appartient à $\mathcal{H}_k$.
        \item \textbf{Propriété de reproduction} : pour toute fonction $f \in \mathcal{H}_k$ et tout $x \in \mathcal{X}$,
        \begin{equation}
            f(x) = \langle f, k_x \rangle_{\mathcal{H}_k}.
        \end{equation}
        \item Pour tous $x, x' \in \mathcal{X}$ :
        \begin{equation}
            k(x, x') = \langle k_x, k_{x'} \rangle_{\mathcal{H}_k} = \langle \varphi_k(x), \varphi_k(x') \rangle_{\mathcal{H}_k} \quad \text{en posant } \varphi_k(x) = k_x.
        \end{equation}
    \end{enumerate}
    L'espace $\mathcal{H}_k$ est appelé l'\emph{espace de Hilbert à noyau reproduisant} (RKHS) associé à $k$.
\end{theorem}

\begin{proof}[Esquisse de preuve]
    La construction s'effectue en trois étapes principales :
    \begin{enumerate}
        \item \textbf{Espace pré-hilbertien} : on pose $\mathcal{H}_0 = \operatorname{vect}\{k_x \mid x \in \mathcal{X}\}$. Pour deux combinaisons linéaires $f = \sum_{i=1}^m a_i k_{x_i}$ et $g = \sum_{j=1}^{m'} b_j k_{z_j}$, on munit $\mathcal{H}_0$ de la forme bilinéaire :
        \begin{equation*}
            \langle f, g \rangle_{\mathcal{H}_0} := \sum_{i=1}^m \sum_{j=1}^{m'} a_i b_j k(x_i, z_j).
        \end{equation*}
        La symétrie et la positivité de $k$ assurent que cette forme est symétrique et positive. De plus, par définition, $\langle f, k_x \rangle_{\mathcal{H}_0} = \sum_{i=1}^m a_i k(x_i, x) = f(x)$.
        
        \item \textbf{Caractère défini} : d'après l'inégalité de Cauchy-Schwarz appliquée à cette forme semi-définie positive, pour tout $x \in \mathcal{X}$ :
        \begin{equation*}
            |f(x)| = |\langle f, k_x \rangle_{\mathcal{H}_0}| \leqslant \|f\|_{\mathcal{H}_0} \sqrt{k(x, x)}.
        \end{equation*}
        Ainsi, $\|f\|_{\mathcal{H}_0} = 0 \implies \forall x \in \mathcal{X}, \; f(x) = 0$, donc $f$ est identiquement nulle. L'espace $\mathcal{H}_0$ est donc un espace pré-hilbertien séparé.
        
        \item \textbf{Complétion fonctionnelle} : on complète $\mathcal{H}_0$ en un espace de Hilbert $\mathcal{H}_k$. Grâce à la majoration ponctuelle issue de Cauchy-Schwarz, toute suite de Cauchy converge ponctuellement vers une fonction bien définie sur $\mathcal{X}$. Les éléments complétés restent donc de véritables fonctions, et la propriété de reproduction s'étend par continuité.
    \end{enumerate}
\end{proof}

Nous avons donc l'équivalence : les noyaux semi-définis positifs sont exactement des produits scalaires de \emph{feature maps}. Choisir un noyau $k$ ou choisir un espace $( \mathcal{H}, \varphi)$ est donc la même chose, sauf que l'on sait explicitement calculer $k$. On peut donc complètement oublier $ \mathcal{H}$ et $ \varphi$ dans le processus d'apprentissage. 

\begin{proposition}[La recette]
    L'algorithme complet d'apprentissage supervisé à noyau se résume donc en 4 étapes : 
    \begin{itemize}
        \item \emph{Choix de $k$} : On définit un noyau $k$ mesurant la "distance" de nos données en fonction de leur forme. 
        \item \emph{Matrice de Gram} : On calcule l'ensemble des paires $K _{i,j} = k(x_i, x_j)$ pour construire la matrice de Gram $K \in \R^{n \times n}$. 
        \item \emph{Optimisation} : On résout un problème d'optimisation pour trouver $ \hat{ \alpha}_S \in \R^n$ : 
        \begin{equation}
            \hat{ \alpha}_S \in \underset{ \alpha \in \R^n}{\arg \min} \frac{1}{n} \sum_{i=1}^{n} l \left(y_i , (K \alpha)_i\right) + \frac{ \lambda}{2} \alpha^\top K \alpha. 
        \end{equation}
        \item \emph{Inférence} : Une fois trouvé, on put prédire le label de tout nouveau $x$ par : 
        \begin{equation}
            \hat{y} = \sum_{j=1}^{n} \hat{ \alpha}_j k(x, x_j). 
        \end{equation}
    \end{itemize}
\end{proposition}

% ======================================================================
% Séparateurs à Vaste Marge (SVM)

\section{Séparateurs à Vaste Marge (SVM)}

L'algorithme des Séparateurs à Vaste Marge (\emph{Support Vector Machines}) illustre parfaitement la puissance des méthodes à noyaux. On commence par formuler le problème géométrique dans le cas linéaire élémentaire, avant d'en dériver la version duale et d'y injecter le \emph{Kernel Trick}.

% ----------------------------------------------------------------------
\subsection{Cas linéairement séparable : Maximisation de la marge}

Soit un échantillon d'apprentissage $S = ((x_1, y_1), \dots, (x_n, y_n)) \in (\R^d \times \{-1, 1\})^n$. On suppose dans un premier temps les données linéairement séparables : il existe un hyperplan passant par l'origine $H(\theta) = \{z \in \R^d \mid \theta^\top z = 0\}$ tel que pour tout $i \in \llbracket 1, n \rrbracket$, $y_i(\theta^\top x_i) > 0$.

\begin{proposition}[Distance à l'hyperplan]
    Pour tout $\theta \neq 0$, la distance orthogonale d'un point $x \in \R^d$ à l'hyperplan $H(\theta)$ est donnée par :
    \begin{equation}
        \operatorname{dist}(x, H(\theta)) = \frac{|\theta^\top x|}{\|\theta\|_2}.
    \end{equation}
\end{proposition}

Sous l'hypothèse de séparabilité, $y_i(\theta^\top x_i) = |\theta^\top x_i|$. L'objectif consiste à trouver l'hyperplan qui maximise la distance minimale aux observations d'entraînement (la marge géométrique) :
\begin{equation}
    \max_{\theta \neq 0} \min_{i \in \llbracket 1, n \rrbracket} \frac{y_i(\theta^\top x_i)}{\|\theta\|_2}.
\end{equation}
Le problème étant invariant par changement d'échelle $\theta \mapsto c\theta$ ($c > 0$), on impose la normalisation canonique $\min_{i} y_i(\theta^\top x_i) = 1$. Maximiser $1/\|\theta\|_2$ sous cette contrainte équivaut à minimiser $\frac{1}{2}\|\theta\|_2^2$. En relâchant l'égalité en une inégalité (qui sera saturée à l'optimum), on obtient le programme quadratique convexe :
\begin{equation}
    \min_{\theta \in \R^d} \; \frac{1}{2} \|\theta\|_2^2 \quad \text{s.c.} \quad \forall i \in \llbracket 1, n \rrbracket, \quad y_i(\theta^\top x_i) \geqslant 1.
\end{equation}

% ----------------------------------------------------------------------
\subsection{Cas non séparable : Variables d'écart et perte charnière}

En pratique, les données se chevauchent ou comportent du bruit, rendant la séparation stricte impossible. Pour autoriser des erreurs tout en les contrôlant, on introduit des variables d'écart (\emph{slack variables}) $\xi_i \geqslant 0$ pénalisées par un coefficient $C > 0$ :
\begin{equation}
    \min_{\theta \in \R^d, \; \xi \in \R^n} \; \frac{1}{2}\|\theta\|_2^2 + C \sum_{i=1}^n \xi_i \quad \text{s.c.} \quad \forall i \in \llbracket 1, n \rrbracket, \quad y_i(\theta^\top x_i) \geqslant 1 - \xi_i \quad \text{et} \quad \xi_i \geqslant 0.
\end{equation}
Ce problème admet toujours des solutions admissibles dès que les $\xi_i$ sont choisis suffisamment grands.

\begin{proposition}[Pont avec la minimisation du risque empirique]
    À $\theta$ fixé, l'objectif est strictement croissant en chaque variable $\xi_i$. À l'optimum, $\xi_i$ atteint donc sa plus petite borne autorisée par les contraintes, soit :
    \begin{equation}
        \xi_i = \max(0, 1 - y_i(\theta^\top x_i)) = \left(1 - y_i(\theta^\top x_i)\right)_+.
    \end{equation}
    En substituant cette expression, en divisant l'objectif par $nC > 0$ et en posant $\lambda := \frac{1}{nC}$, le problème géométrique coïncide rigoureusement avec un problème d'ERM régularisé pour la perte charnière (\emph{hinge loss}) :
    \begin{equation}
        \min_{\theta \in \R^d} \; \frac{1}{n} \sum_{i=1}^n \left(1 - y_i(\theta^\top x_i)\right)_+ + \frac{\lambda}{2}\|\theta\|_2^2.
    \end{equation}    
\end{proposition}


% ----------------------------------------------------------------------
\subsection{Dualité de Wolfe et formulation duale}

Le problème primal étant convexe avec des contraintes affines, les conditions de qualification de Slater sont vérifiées (par exemple pour $\theta = 0$ et $\xi_i = 2$), garantissant la dualité forte.

On introduit les multiplicateurs de Lagrange $\alpha_i \geqslant 0$ (pour les contraintes de marge) et $\beta_i \geqslant 0$ (pour la positivité de $\xi_i$) :
\begin{equation}
    \mathcal{L}(\theta, \xi, \alpha, \beta) = \frac{1}{2}\|\theta\|_2^2 + C \sum_{i=1}^n \xi_i - \sum_{i=1}^n \alpha_i \left(y_i \theta^\top x_i + \xi_i - 1\right) - \sum_{i=1}^n \beta_i \xi_i.
\end{equation}

\paragraph{}
\begin{proposition}
    Minimisons le Lagrangien : 
    \begin{itemize}
        \item \textbf{En $\xi$ :} Le Lagrangien s'écrit $\sum_{i=1}^n (C - \alpha_i - \beta_i)\xi_i + \dots$ L'infimum par rapport à $\xi_i \geqslant 0$ vaut $-\infty$ sauf si $C - \alpha_i - \beta_i = 0$. Cette condition permet d'éliminer $\beta_i \geqslant 0$, ce qui équivaut à la contrainte de boîte :
        \begin{equation}
            0 \leqslant \alpha_i \leqslant C \quad \forall i \in \llbracket 1, n \rrbracket.
        \end{equation}
        \item \textbf{En $\theta$ :} La condition d'annulation du gradient $\nabla_\theta \mathcal{L} = \theta - \sum_{i=1}^n \alpha_i y_i x_i = 0$ impose :
        \begin{equation}
            \theta^* = \sum_{i=1}^n \alpha_i y_i x_i.
        \end{equation}
    \end{itemize}
\end{proposition}

En réinjectant ces relations dans $\mathcal{L}$, le problème dual de Wolfe s'écrit sous forme d'un programme quadratique concave :
\begin{equation}
    \max_{\alpha \in \R^n} \; \sum_{i=1}^n \alpha_i - \frac{1}{2} \sum_{i=1}^n \sum_{j=1}^n \alpha_i \alpha_j y_i y_j (x_i^\top x_j) \quad \text{s.c.} \quad \forall i \in \llbracket 1, n \rrbracket, \quad 0 \leqslant \alpha_i \leqslant C.
\end{equation}

\begin{remark}
    Les données d'entrée $x_i$ n'interviennent dans ce problème dual que par leurs produits scalaires $x_i^\top x_j$.
\end{remark}

% ----------------------------------------------------------------------
\subsection{Le Kernel SVM (SVM à noyau)}

Soit $k$ un noyau semi-défini positif sur $\mathcal{X}$. En vertu du théorème d'Aronszajn, $k$ définit implicitement un RKHS $\mathcal{H}_k$ et une application canonique $\varphi_k$ tels que $k(x, x') = \langle \varphi_k(x), \varphi_k(x') \rangle_{\mathcal{H}_k}$. 
On formule le problème primal directement dans $\mathcal{H}_k$ :
\begin{equation}
    \min_{\theta \in \mathcal{H}_k} \; \frac{1}{n} \sum_{i=1}^n \left(1 - y_i \langle \theta, \varphi_k(x_i) \rangle_{\mathcal{H}_k}\right)_+ + \frac{\lambda}{2} \|\theta\|_{\mathcal{H}_k}^2.
\end{equation}

Grâce au \emph{Kernel Trick}, on substitue directement chaque produit scalaire par l'évaluation du noyau $k(x_i, x_j) = K_{i,j}$ dans la formulation duale. En posant le changement de variable $\mu_i = \lambda \alpha_i$ (avec $\lambda = \frac{1}{nC}$, ce qui donne $0 \leqslant \mu_i \leqslant \frac{1}{n}$), le problème dual s'écrit sous forme matricielle :
\begin{equation}
    \boxed{\max_{\mu \in \R^n} \; \sum_{i=1}^n \mu_i - \frac{1}{2\lambda} \sum_{i=1}^n \sum_{j=1}^n y_i y_j \mu_i \mu_j k(x_i, x_j) \quad \text{s.c.} \quad \forall i \in \llbracket 1, n \rrbracket, \quad 0 \leqslant \mu_i \leqslant \frac{1}{n}.}
\end{equation}
Il s'agit d'un problème d'optimisation quadratique à $n$ variables, dont la résolution numérique s'effectue uniquement à l'aide de la matrice de Gram $K \in \mathcal{S}_n^+(\R)$.

\begin{proposition}[Règle de décision et vecteurs de support]
    D'après l'expression de $\theta^*$ et la propriété de reproduction dans le RKHS, la fonction de décision optimale pour classer un nouveau point $x \in \mathcal{X}$ est :
    \begin{equation}
        f(x) = \operatorname{sign}\left( \sum_{i=1}^n \hat{\alpha}_i y_i k(x_i, x) \right).
    \end{equation}
    D'après les conditions de complémentarité de Karush-Kuhn-Tucker (KKT), les coefficients $\hat{\alpha}_i$ sont strictement nuls pour tous les points situés en dehors de la marge. La décision ne dépend donc que des points vérifiant $\hat{\alpha}_i > 0$, appelés \emph{vecteurs de support} :
    \begin{equation}
        f(x) = \operatorname{sign}\left( \sum_{i \in \mathrm{Support}} \hat{\alpha}_i y_i k(x_i, x) \right).
    \end{equation}
    Ce résultat garantit la parcimonie du modèle et rend l'évaluation en phase d'inférence très rapide.
\end{proposition}
